% VI24_FULL_SOLUTION_REFINEMENT_V1
\subsection*{VI.24.P2: Legacy CA13-E7: scalar extension of a polynomial quotient}
\begin{problem}\label{prob:vi24-p02}
Let $k\subseteq K$ be a field extension and $f\in k[x]$. Prove
\[
(k[x]/(f))\otimes_kK\cong K[x]/(f).
\]
\end{problem}

\ifdefined\FullProblemDossiers
\paragraph{Why this problem matters.}
This dossier isolates a reusable piece of the base-change calculus and records both its categorical and affine content.

\paragraph{Useful ideas before starting.}
Start from the Cartesian universal property.  When the schemes are affine, translate the square contravariantly into tensor products, quotients, or localizations, and then translate the resulting ring statement back into geometry.

\paragraph{Guided hints.}
\begin{enumerate}
\item Identify the relevant pullback or scalar-extension ring.
\item Use the appropriate universal property or exactness statement.
\item Check that the canonical maps are the expected projections.
\item State the geometric consequence, not only the ring isomorphism.
\end{enumerate}

\begin{solution}
Begin with the exact sequence of $k$-vector spaces
\[
(f)\longrightarrow k[x]\longrightarrow k[x]/(f)\longrightarrow0.
\]
Tensor over $k$ with $K$.  Since $K$ is a flat $k$-module (indeed every vector space over a field is flat), exactness is preserved, and Exercise VI.24.3 gives
\[
k[x]\otimes_kK\cong K[x].
\]
The image of $(f)\otimes_kK$ in $K[x]$ is the extended ideal generated by the same polynomial $f$.  Hence the cokernel is
\[
(k[x]/(f))\otimes_kK\cong K[x]/(f).
\]
The isomorphism sends $(g(x)\bmod f)\otimes a$ to $a g(x)\bmod f$.  Geometrically, the hypersurface cut out by $f$ is transported to the new field by keeping the equation and extending its coefficients.  What can change is the factorization of $f$, and therefore the component structure of the scalar-extended scheme.
\end{solution}

\paragraph{What happened mathematically.}
The base-changed object is forced by a universal property; tensor products make that universal property computable on affine charts.

\paragraph{Extensions.}
The same pattern will be reused in VI/25 for fibers and in VI/26 for properties of morphisms under change of base.

\paragraph{Provenance.}
\textbf{Legacy CA13-E7.} This reserved scalar-extension exercise belongs to VI/24, not VI/23.
\else
\paragraph{Provenance.} \textbf{Legacy CA13-E7.} This reserved scalar-extension exercise belongs to VI/24, not VI/23.
\fi
